\documentclass{scrreprt}
\usepackage{listings}
\usepackage{underscore}
\usepackage{graphicx}
\usepackage[bookmarks=true]{hyperref}
\usepackage[utf8]{inputenc}
\usepackage[english]{babel}
\usepackage{hyperref}
\usepackage{tabularx}
\usepackage{imakeidx}
\makeindex[intoc]
\usepackage{graphicx}
\graphicspath{ {./Pictures/} }

\hypersetup{
    bookmarks=false,    % show bookmarks bar?
    pdftitle={Software Requirement Specification},    % title
    pdfauthor={Sara Lindholm},                     % author
    pdfsubject={TeX and LaTeX},                        % subject of the document
    pdfkeywords={TeX, LaTeX, graphics, images}, % list of keywords
    colorlinks=true,       % false: boxed links; true: colored links
    linkcolor=blue,       % color of internal links
    citecolor=black,       % color of links to bibliography
    filecolor=black,        % color of file links
    urlcolor=purple,        % color of external links
    linktoc=page            % only page is linked
}%
\def\myversion{1.3}
\title{Software Requirements Specification}
\author{Sara Lindholm \\ Patrick Asman \\ Jacob Rhan}

\date{October 2020}

\newcolumntype{s}{>{\hsize=3.75\hsize}X}
\newcolumntype{m}{>{\hsize=7\hsize}X}
\newcolumntype{b}{>{\hsize=11\hsize}X}

\begin{document}

\begin{titlepage}
    \begin{center}
    \begin{bfseries}
        \Huge{SOFTWARE REQUIREMENTS SPECIFICATION}\\
        \vspace{1.5cm}
        \LARGE Analysis Team \\
        \vspace{1.5cm}
        %Submitted to : - \\Lecturer\\
        %\vspace{1.5cm}
        \today\\
        %\vspace{5cm}
        \vspace{1.5cm}
        {Version \myversion}\\
        \vfill
        \includegraphics[width=\linewidth]{Pictures/logo.png} \\
    \end{bfseries}
    \end{center}
\end{titlepage}

\tableofcontents

\chapter*{Revision History}
\begin{center}
\begin{tabular}{|p{1,9cm}|p{1,5cm}|p{6cm}|p{4cm}|}
 \hline
 \textbf{Date} & \textbf{Version} & \textbf{Description} & \textbf{Author} \\ 
 \hline
 2020-10-11 & 0.9 & Initial SRS structure & Sara Lindholm \\
 \hline
 2020-10-15 & 1.0 & First Draft & Patrick Asman, Sara Lindholm, Jacob Rhan \\ 
 \hline
 2020-11-11 & 1.1 & Updated and added functional requirements & Sara Lindholm \\
 \hline
 2020-11-17 & 1.2 & Added category descriptions & Jacob Rhan, Sara Lindholm\\
 \hline
 2020-11-17 & 1.3 & Revised and updated functional requirements with titles and preconditions, as well as updated table format for non-functional requirements. & Patrick Asman, Sara Lindholm, Jacob Rhan, Hjalmar Svensson \\
 \hline
\end{tabular}
\end{center}



\chapter{Introduction}
This section gives an overview of everything included in this document. The purpose and scope for this document is described and a list of abbreviations, acronyms and definitions is provided. 

\section{Purpose}
The purpose of this document is to give a detailed description of \textcolor{red}{[system name?]}. It will illustrate the purpose and features of the system. It will also explain system constraints, interfaces of the system and interactions with external applications. This document is intended primarily to be proposed to the stakeholders for approval and as a reference for the developers.

\section{Scope}
\textcolor{red}{[System name?]} is a web-based application to aid Region Östergötland's health care staff in overseeing and handling many patients. The system will make it easier for the health care staff to oversee many patients' measured health data and prioritize which patients are in most need of care. The system will have customizable overviews, customizable alert-levels and sorting and filtering functions to make prioritization easier. The system will in addition to a better overview of all patients be able to give more detailed overview about the individual patients' health information, such as current medications, health conditions and the measured data. 


\section{Definitions, acronyms and abbreviations}
\begin{center}
\begin{tabularx}{\linewidth}{| l | X |}
 \hline
 \textbf{Term} & \textbf{Definition} \\ 
 \hline
  User  & Health care staff who interacts with the web application.  \\
  \hline
  Admin/Administrator & System administrator who is given specific permission for managing and controlling the system  \\ 
  \hline
  RÖ  & Region Östergötland  \\ 
  \hline
   "Verksamhet" & \\ 
  \hline 
  "Enhet" & \\
  \hline
  Team & \\ 
  \hline
  User group & \\ 
  \hline
\end{tabularx}
\end{center}

\section{References}
IEEE Recommended Practice for Software Requirements Specifications," in IEEE Std 830-1998 , vol., no., pp.1-40, 20 Oct. 1998, doi: 10.1109/IEEESTD.1998.88286.

\section{Overview}
The second chapter of the document, Overall Description, gives an overview of the functionality of the product.


The third chapter, Requirements Specification, gives a more  detailed description of the system and its features. 

\chapter{Overall Description}
This section will present an overview of the entire system. The system's basic functionality and how it interacts with other systems will be introduced. % The constraints and assumptions for the system will also be presented. 

\section{Product Perspective}
\textbf{TODO:} fix. \\ put the product into perspective with other related products
\begin{itemize}
    \item web application
    \item communicates with openEHR (store and recieve health data?), OAuth, RÖ's personell lists, Swedish population register and online health care services (such as 1177, Fass etc).
\end{itemize}


\section{Product functions}
This web application will aid health care staff in handling many patients and their health data. The system will have an overview in which the user will be able to see all patients in their care unit. These patients can then be sorted and filtered by different criteria, which will aid the user in prioritizing which patients are in most need of care. 


The user will also be able to show more detailed information about the individual patients should they choose to click on a patient. In the individual patient view the user will get access to more information than what was available in the overview of all patients. The individual patient may have measured health data, medical information and history, and scheduled appointments which will be displayed in different tabs on the web application.


The application will also be able to send out notifications to the users when there are updates on patients. The specifics for notifications are to be set by the admin so as to avoid being overwhelmed by notifications. 

\section{User Characteristics}
There are two intended users of the system: users and administrators. These two types of users will have different roles and therefore not have access to all the same functionalities. Some of the customizable elements shall only be customizable by an admin and the admin will also be the one to set the criteria for the different alert levels as well as the notifications.


\section{Constraints}


\section{Assumptions and dependencies}
%"This subsection of the SRS should list each of the factors that affect the requirements stated in the SRS. These factors are not design constraints on the software but are, rather, any changes to them that can affect the requirements in the SRS. For example, an assumption may be that a specific operating system will be available on the hardware designated for the software product. If, in fact, the operating system is not available, the SRS would then have to change accordingly."

\section{Apporting of requirements}
%There are some requirements that could be delayed until future versions of the systems if the project is delayed.

\chapter{Specific requirements}
This chapter contains the software requirements and gives a detailed
description of the system and its features. 

\section{External interface requirements}
\textbf{TODO: fix}
\begin{itemize}
    \item openEHR
    \item OAuth
    \item Region Östergötland's personell lists 
    \item Swedish population register 
    \item online health care services (e.g. rådgivningsstöd, 1177, vårdhandboken, FASS, internetmedicin) 
\end{itemize}

%\subsection{User interfaces}
%\subsection{Software interfaces}
%\subsection{Communication interfaces}

\section{Functional requirements}
The following section includes the requirements that specify all the fundamental actions of the software system. The section is divided into four main categories: "Admin", "Patient", "System" and "Översikt". These categories are then divided into subcategories based on the different views in each category. 
\\ \\
To clarify how the requirements presented in tables see explanation below. 

\begin{center}
\begin{tabularx}{\linewidth}{| s | m | m | b | s |}
\hline
\textbf{ID:} & \textbf{Title} & \textbf{Precondition} & \textbf{Description} & \textbf{Priority} \\
\hline
FRXXX id of requirement & 
Title to briefly explain the requirement & 
Explains any preconditions for requirement & 
Description of what will happen in the system.  
    \newline \textbf{Input:} What triggers the function  
    \newline \textbf{Action:} How the system responds to the input 
    \newline \textbf{Output:} The function's output &
Priority level (Low-Medium-High) \\ 
\hline
\end{tabularx}
\end{center}

\subsection{Admin}
In section 3.2.1. we describe functional requirements for the Admin. A precondition is that the admin is logged in to the system. 
\subsubsection{Översikt/Notislogg}
A precondition is that the admin is logged in and on the view “Oversikt/Notislogg”
\begin{center}
\begin{tabularx}{\linewidth}{| s | m | m | b | s |}
\hline
\textbf{ID:} & \textbf{Title} & \textbf{Precondition} & \textbf{Description} & \textbf{Priority} \\
\hline
FR001 & 
Customization of views &
New-view-button has been pressed and system currently displays "Skapa ny överblick" window. & 
\textbf{Input:} Filter settings selected and "OK"-button pressed.  
    \newline \textbf{Action:} System generates new view and saves to database.
    \newline \textbf{Output:} New view available for all users in the selected user group. &
Medium \\ 
\hline
FR002 & 
Customization of overview dashboard &
& 
\textbf{Input:} Display-settings selected and ""OK""-button pressed.
    \newline \textbf{Action:} System generates and saves display-settings. 
    \newline \textbf{Output:} Overview dashboard updated according to selected settings. & 
Medium \\
\hline
\end{tabularx}
\end{center}

\subsubsection{Översikt/Patienter}
A precondition is that the admin is logged in and on the view “Oversikt/Patienter”
\begin{center}
\begin{tabularx}{\linewidth}{| s | m | m | b | s |}
\hline
\textbf{ID:} & \textbf{Title} & \textbf{Precondition} & \textbf{Description} & \textbf{Priority} \\
\hline
FR003 & 
Customization of views &
New-view-button has been pressed and system currently displays "Skapa ny överblick" window. & 
\textbf{Input:} Filter settings selected and "OK"-button pressed.  
    \newline \textbf{Action:} System generates new view and saves to database.
    \newline \textbf{Output:} New view available for all users in the selected user group. &
Medium \\ 
\hline
FR004 & 
Customization of overview dashboard &
& 
\textbf{Input:} Display-settings selected and ""OK""-button pressed.
    \newline \textbf{Action:} System generates and saves display-settings. 
    \newline \textbf{Output:} Overview dashboard updated according to selected settings. & 
Medium \\
\hline
\end{tabularx}
\end{center}

\subsection{Patient}
In section 3.2.2. we describe functional requirements for the User. A precondition is that the user is logged in to the system and, on the view “Patient”

\begin{center}
\begin{tabularx}{\linewidth}{| s | m | m | b | s |}
\hline
\textbf{ID:} & \textbf{Title} & \textbf{Precondition} & \textbf{Description} & \textbf{Priority} \\
\hline
FR005 & 
Switch to overview view & 
 & 
    \newline \textbf{Input:} User clicks on the Home button.
    \newline \textbf{Action:} Loading Overview view.
    \newline \textbf{Output:} System displays Overview View. &
High \\ 
\hline
\end{tabularx}
\end{center}

\subsubsection{Admin}
A precondition is that the admin is logged in and on the view “Patient/Admin” 

\begin{center}
\begin{tabularx}{\linewidth}{| s | m | m | b | s |}
\hline
\textbf{ID:} & \textbf{Title} & \textbf{Precondition} & \textbf{Description} & \textbf{Priority} \\
\hline
FR006 & 
Set limit values for priority levels & 
 & 
    \textbf{Input}: Limit values for different measurements and click save
    \newline \textbf{Action}: System saves and sets patient specific prioritization limits
    \newline \textbf{Output}: Settings displayed &
High \\ 
\hline 
FR007 & 
Accessing activity log & 
 & 
\textbf{Input}: Click on ""Admin"" tab \newline
\textbf{Action}: Load activity log from database \newline 
\textbf{Output}: Display activity log &
Medium \\ 
\hline 
FR008 & 
Search in activity log & 
 & 
\textbf{Input}: Text string \newline 
\textbf{Action}: System checks for activities or users that matches text string \newline 
\textbf{Output}: Displays activities and/or users that matches string. &
Low \\ 
\hline
FR009 & 
Filter in activity log & 
Filter-button has been pressed and system currently displays "Filter" window.  & 
\textbf{Input}: Filter settings selected and "OK"-button pressed. \newline 
\textbf{Action}: System filter out activites from filter settings. \newline 
\textbf{Output}: Displays activities that matches filter." &
Low \\ 
\hline
\end{tabularx}
\end{center}

\subsubsection{Kalender}
A precondition is that the user is logged in and on the view “Patient” and, on the tab “Kalender”
\begin{center}
\begin{tabularx}{\linewidth}{| s | m | m | b | s |}
\hline
\textbf{ID:} & \textbf{Title} & \textbf{Precondition} & \textbf{Description} & \textbf{Priority} \\
\hline
FR010 & 
Add calendar activity &
& 
\textbf{Input}: Click "Add Activity"-button \newline 
\textbf{Action}: Activity added \newline 
\textbf{Output}: Display the new activity & 
Low \\ 
\hline
\end{tabularx}
\end{center}

\subsubsection{Mätvärden}
A precondition is that the user is logged in and on the view “Patient” and, on the tab “Mätvärden”
\begin{center}
\begin{tabularx}{\linewidth}{| s | m | m | b | s |}
\hline
\textbf{ID:} & \textbf{Title} & \textbf{Precondition} & \textbf{Description} & \textbf{Priority} \\
\hline
FR011 & 
Detailed graph and list of measurement &
& 
\textbf{Input}: Click on a graph \newline 
\textbf{Action}: Load diagram and list of measurements \newline 
\textbf{Output}: Display diagrams and list of measurements & 
Low \\ 
\hline 
FR012 & 
Update patient's measurement values &
& 
\textbf{Input}: Fill in the updated measurements and click "Signera" \newline 
\textbf{Action}: System store measured value \newline 
\textbf{Output}: Displays confirmation message & 
High \\ 
\hline 
FR013 & 
Accept patient's measurement values &
User has clicked on the notification symbol and system currently displays the "Uppmärksammat värde" window. & 
\textbf{Input}: Accept unexpected measurement value by clicking "Signera" \newline 
\textbf{Action}: System store measured value \newline 
\textbf{Output}: Displays changes & 
Medium \\ 
\hline 
FR014 & 
Deny patient's measurement values &
User has clicked on the notification symbol and system currently displays the "Uppmärksammat värde" window. & 
\textbf{Input}: Deny unexpected measurement value by clicking "Avbryt" \newline 
\textbf{Action}: System remove measured \newline 
\textbf{Output}: Displays changes & 
Medium \\ 
\hline 
\end{tabularx}
\end{center}

\subsection{System}
In section 3.2.3. we describe functional requirements for the system. A precondition is that the user is logged in to the system.
\begin{center}
\begin{tabularx}{\linewidth}{| s | m | m | b | s |}
\hline
\textbf{ID:} & \textbf{Title} & \textbf{Precondition} & \textbf{Description} & \textbf{Priority} \\
\hline
FR015 & 
Log users accessing patient information &
& 
\textbf{Input}: Patient information is accessed by user \newline 
\textbf{Action}: System logs user ID, save to database and updates activity log. \newline 
\textbf{Output}: Updated activity log. & 
High \\ 
\hline 
FR016 & 
Log users updating patient information &
& 
\textbf{Input}: Patient information (e.g measurements) is updated by user \newline 
\textbf{Action}: System logs user ID, saves to database and updates activity log. \newline 
\textbf{Output}: Updated activity log. &
High \\ 
\hline 
FR017 & 
Not submitted data notification &
& 
\textbf{Input}: Time limit, according to individual health care plan, expires \newline 
\textbf{Action}: Notification printing and sending to responsible user \newline 
\textbf{Output}: The system sends notification to user & 
High \\ 
\hline 
FR018 & 
Limit notification &
FR006 is fulfilled & 
\textbf{Input}: Patient measurement \newline 
\textbf{Action}: Compare to set limit value \newline 
\textbf{Output}: Send notification & 
High \\ 
\hline 
FR019 & 
Log out & 
&
\textbf{Input}: User clicks on the log out button. \newline 
\textbf{Action}: System logs out user and loads start page. \newline 
\textbf{Output}: System displays start page. & 
High \\ 
\hline 
\end{tabularx}

\begin{tabularx}{\linewidth}{| s | m | m | b | s |}
\hline
FR020 & 
System help &
 & 
\textbf{Input}: User clicks on button "Help" \newline 
\textbf{Action}: Loading help information message \newline 
\textbf{Output}: The system displays help information message & 
Low \\ 
\hline 
FR021 & 
Patient contact alternatives &
 & 
\textbf{Input}: Patient identification and contact information \newline
\textbf{Action}: The system loads contact information about patient \newline 
\textbf{Output}: Display patient's contact information as notification & 
Low \\ 
\hline 
\end{tabularx}
\end{center}

\subsection{Översikt}
In section 3.2.4. we describe functional requirements for the User. A precondition is that the user is logged in to the system and, on the view “Översikt”
he system confirms that user wants to access patient data not in their department.

\subsubsection{Kalender}
A precondition is that the user is logged in, on the view “Översikt” and, on the tab “Kalender”
\begin{center}
\begin{tabularx}{\linewidth}{| s | m | m | b | s |}
\hline
\textbf{ID:} & \textbf{Title} & \textbf{Precondition} & \textbf{Description} & \textbf{Priority} \\
\hline
FR022 & 
Add calendar activity & 
&
\textbf{Input}: Click "Add Activity"-button \newline 
\textbf{Action}: Activity added \newline
\textbf{Output}: Display the new activity & 
Low \\ 
\hline
\end{tabularx}
\end{center}

\subsubsection{Notislogg}
A precondition is that the user is logged in, on the view “Översikt” and, on the tab “Notislogg”

\begin{center}
\begin{tabularx}{\linewidth}{| s | m | m | b | s |}
\hline
\textbf{ID:} & \textbf{Title} & \textbf{Precondition} & \textbf{Description} & \textbf{Priority} \\
\hline
FR023 & 
Active choice of accessing patient information outside of user's unit & 
User in a certain unit tries to access patient information about a patient not in user's unit. &
\textbf{Input}: Clicking on a patient in overview \newline
\textbf{Action}: The system checks if user wants to access patient data. \newline
\textbf{Output}: Display message checking if user wants to access data. & 
High \\ 
\hline
FR024 & 
Accessing patient information outside of user's unit & 
User in a certain unit tries to access patient information about a patient not in user's unit. &
\textbf{Input}: Clicking on "Bekräfta" button to access patient data \newline
\textbf{Action}: The system confirms that data shall be displayed \newline 
\textbf{Output}: The system displays patient data & 
High \\ 
\hline
FR025 & 
Sorting patients & 
&
\textbf{Input}: sorting values from categories [age, priority-level, measurements, health care unit, team, disease, patients’conditions, disease type, latest date, name, updated by]. \newline 
\textbf{Action}: The system rearranges the list of patients according to sorting criteria. \newline 
\textbf{Output}: The system displays all patients according to the sorting criteria. & 
High \\ 
\hline
FR026 & 
Filtering patients & 
&
\textbf{Input}: filter values from categories [age, alert-level, measurements, health care unit, team, disease and patients’conditions among others]. \newline
\textbf{Action}: The system searches for patients matching filter criteria. \newline
\textbf{Output}: The system displays all patients who matched the filters. & 
High \\ 
\hline
\end{tabularx}

\begin{tabularx}{\linewidth}{| s | m | m | b | s |}
\hline
FR027 & 
Switch to patient view & 
&
\textbf{Input}: User clicks on a patient. \newline
\textbf{Action}: Loading Patient view. \newline
\textbf{Output}: System displays Patient View. & 
High \\ 
\hline
FR028 & 
Hover over notification to present detailed information  & 
Active hovering over a notification in "Notislogg" &
\textbf{Input}: Mouse hover over notification \newline
\textbf{Action}: System loads data \newline
\textbf{Output}: Displays detailed information & 
Medium \\ 
\hline
FR029 & 
Dynamic diagnosis filtering & 
&
\textbf{Input}: Text string in diagnosis filter \newline 
\textbf{Action}: System checks for diagnosis that matches text string \newline
\textbf{Output}: Displays diagnosis that matches string dynamically & 
Medium \\ 
\hline
FR030 & 
Visible sorting parameters & 
&
\textbf{Input}: Click sort \newline 
\textbf{Action}: Data loading \newline
\textbf{Output}: Information displayed in relation to search result & 
Low \\ 
\hline
\end{tabularx}
\end{center}

\subsubsection{Patienter}
A precondition is that the user is logged in, on the view “Översikt” and, on the tab “Patienter”

\begin{center}
\begin{tabularx}{\linewidth}{| s | m | m | b | s |}
\hline
\textbf{ID:} & \textbf{Title} & \textbf{Precondition} & \textbf{Description} & \textbf{Priority} \\
\hline
FR031 & 
Filtering patients & 
&
\textbf{Input}: filter values from categories [age, alert-level, measurements, health care unit, team, disease and patients’conditions among others]. \newline
\textbf{Action}: The system searches for patients matching filter criteria. \newline
\textbf{Output}: The system displays all patients who matched the filters. & 
High \\ 
\hline
FR032 & 
Active choice of accessing patient information outside of user's unit & 
User in a certain unit tries to access patient information about a patient not in user's unit. &
\textbf{Input}: Clicking on a patient in overview \newline
\textbf{Action}: The system checks if user wants to access patient data. \newline
\textbf{Output}: Display message checking if user wants to access data. & 
High \\ 
\hline
FR033 & 
Accessing patient information outside of user's unit & 
User in a certain unit tries to access patient information about a patient not in user's unit. &
\textbf{Input}: Clicking on "Bekräfta" button to access patient data \newline
\textbf{Action}: The system confirms that data shall be displayed \newline 
\textbf{Output}: The system displays patient data & 
High \\ 
\hline
FR034 & 
Sorting patients & 
&
\textbf{Input}: sorting values from categories [age, priority-level, measurements, health care unit, team, disease, patients’conditions, disease type, latest date, name, updated by]. \newline 
\textbf{Action}: The system rearranges the list of patients according to sorting criteria. \newline 
\textbf{Output}: The system displays all patients according to the sorting criteria. & 
High \\
\hline
\end{tabularx}

\begin{tabularx}{\linewidth}{| s | m | m | b | s |}
\hline
FR035 & 
Switch to patient view & 
&
\textbf{Input}: User clicks on a patient. \newline
\textbf{Action}: Loading Patient view. \newline
\textbf{Output}: System displays Patient View. & 
High \\ 
\hline
FR036 & 
Dynamic diagnosis filtering & 
&
\textbf{Input}: Text string in diagnosis filter \newline 
\textbf{Action}: System checks for diagnosis that matches text string \newline
\textbf{Output}: Displays diagnosis that matches string dynamically & 
Medium \\ 
\hline
FR037 & 
Hover over notification to present detailed information  & 
Active hovering over a notification in "Notislogg" &
\textbf{Input}: Mouse hover over notification \newline
\textbf{Action}: System loads data \newline
\textbf{Output}: Displays detailed information & 
Medium \\
\hline
\end{tabularx}
\end{center}

\section{Performance requirements}
The requirements in this section provide specification of the system performance. 
\begin{center}
\begin{tabularx}{\linewidth}{| l | X | l |}
\hline
\textbf{ID:} & \textbf{Statement} & \textbf{Priority} \\
\hline
NFR005 & 
The system shall be able to store and handle 10 000 simultaneous active patients. & 
High \\ 
\hline
\end{tabularx}
\end{center}

\section{Software system attributes}
The requirements in this section specify the required usability, security, maintainability and portability.
of the software system.

\subsection{Usability}

\begin{center}
\begin{tabularx}{\linewidth}{| l | X | l |}
\hline
\textbf{ID:} & \textbf{Statement} & \textbf{Priority} \\
\hline
NFR008 & 
90\% of the system users shall perceive the system as easy to use and flexible. & 
Low \\ 
\hline
NFR012 & 
The Google Material Design system themed with Region Östergötland’s color palette shall be used. & 
Medium \\ 
\hline
NFR016 & 
The system shall be easy to navigate (e.g. low Lostness-score, follow the ISO 9241-11 standard). &
Low \\ 
\hline
NFR017 & 
The system shall provide good user experience (follow the ISO 9241-11 standard). & 
Low \\ 
\hline
NFR018 & 
The system shall be accessible for people with disabilities and follow WCAG guidelines. &
Low \\ 
\hline
\end{tabularx}
\end{center}

\subsection{Security}

\begin{center}
\begin{tabularx}{\linewidth}{| l | X | l |}
\hline
\textbf{ID:} & \textbf{Statement} & \textbf{Priority} \\
\hline
NFR002 & 
OAuth shall be used for login and authentication. &
High \\ 
 \hline
\end{tabularx}
\end{center}

\subsection{Maintainability}

\begin{center}
\begin{tabularx}{\linewidth}{| l | X | l |}
\hline
\textbf{ID:} & \textbf{Statement} & \textbf{Priority} \\
\hline
NFR013 & 
The system shall enable use of openEHR platforms as storage backend for EHR content (e.g. measurements, content from forms filled out by the patient and care plan) &
High \\ 
 \hline
NFR014 & 
The system shall be integrable with RÖ’s systems for personell lists from their database.&
High \\ 
\hline 
NFR015 & 
The system shall be integrable with the Swedish population register (folkbokföringsregister).& 
High \\ 
\hline
NFR019 & 
Online health care services (i.e. rådgivningsstöd, 1177, vårdhandboken, FASS, internetmedicin) shall be integrated in the system.& 
Low \\ 
\hline
\end{tabularx}
\end{center}

\subsection{Portability}

\begin{center}
\begin{tabularx}{\linewidth}{| l | X | l |}
\hline
\textbf{ID:} & \textbf{Statement} & \textbf{Priority} \\
\hline
NFR001 & 
The system shall be OS independent. &
High \\ 
\hline
NFR004 & 
It shall be possible to deploy the system using Docker orchestrated by Kubernetes.&
High \\ 
\hline
NFR006 & 
Standardized terms and national language for special purposes (LSP) shall be used. &
Low \\ 
\hline
NFR007 & 
The system shall be written in TypeScript (instead of JavaScript).&
Low \\ 
\hline
NFR009 & 
The system shall implement a scalable UI that works seamlessly across a broad range of screen sizes. &
Medium \\ 
\hline
\end{tabularx}
\end{center}


\section{Other requirements}

\begin{center}
\begin{tabularx}{\linewidth}{| l | X | l |}
\hline
\textbf{ID:} & \textbf{Statement} & \textbf{Priority} \\
\hline
NFR003 & 
The server-side functionality shall be based on open source parts. &
High \\ 
\hline
NFR010 & 
The system shall be developed in open source, Apache 2 license. &
Medium \\ 
\hline
NFR011 & 
The graphical interface components shall be continuously and openly documented/tested/published through Storybook. &
Medium \\ 
\hline
\end{tabularx}
\end{center}


\textbf{TODO}: add and categorize non-functional requirements (NFR020-NFR036) from excel-file v. 1.7 


\appendix
\chapter{Appendix Title}

%\printindex

\end{document}